\section{Related Work}
\subsection{CLIP-based Prompt Adaptation for Anomaly Detection and Segmentation}
Most zero-shot industrial AD adapts CLIP~\cite{radford2021}, despite it not being designed for inspection. Prompt-based methods craft or learn normal/abnormal prompts: WinCLIP~\cite{jeong2023winclip} (window aggregation), AnomalyCLIP~\cite{zhou2024anomalyclip} (object-agnostic prompts), and PromptAD~\cite{Li_2024_WACV}/AdaCLIP~\cite{Cao_2024} (learnable prompt tokens). For localization, APRIL-GAN~\cite{chen2023zero}, AnoCLIP~\cite{deng2024}, and SimCLIP~\cite{simclip} improve patch-level visual--text alignment; CLIP-SAM~\cite{li2024clipsamclipsamcollaboration} and FiLo~\cite{gu2024filozeroshotanomalydetection} add external segmentation/grounding modules; and MultiADS~\cite{sadikaj2025multiadsdefectawaresupervisionmultitype} adds defect-aware supervision for multi-type segmentation. Others rely on patch comparison or reconstruction~\cite{li2024musczeroshotindustrialanomaly,schwartz2024maeday,aota2023zero}. These methods largely depend on CLIP-style embedding comparison or external modules, making them less suited to unified defect reasoning and analysis.

\subsection{Autoregressive VLMs for Anomaly Detection and Reasoning}
Autoregressive VLMs offer a different paradigm, enabling open-ended visual reasoning and instruction following instead of fixed-prompt embedding comparison. AnomalyGPT~\cite{gu2023anomalygptdetectingindustrialanomalies} is an early effort, but it relies on an external CLIP-based anomaly map combined with the input image and fed to the LLM, so the VLM mainly serves as a reasoning interface over a separately generated map. More recently, Anomaly-OV~\cite{xu2025towards} specializes MLLMs for zero-shot anomaly detection and reasoning via the Anomaly-Instruct-125k instruction set, but targets image-level detection and explanation, with pixel-level segmentation not grounded in the decoding process.

In contrast, GenAU treats the autoregressive VLM as the core anomaly understanding model rather than as a reasoning interface over externally generated anomaly maps. Building on the reasoning-segmentation paradigm of LISA~\cite{lai2024lisareasoningsegmentationlarge}, in which a VLM emits a segmentation token whose hidden state drives a mask, GenAU uses decoder hidden states of dedicated \texttt{[SEG\_normal]} and \texttt{[SEG\_defect]} tokens as language-grounded anchors, aligning them directly with frozen multi-scale visual features (without an external mask decoder) to generate pixel-level anomaly maps. This design unifies image-level anomaly detection, anomaly segmentation, multi-type anomaly detection, and textual defect analysis in a single instruction-following framework. An extended discussion is provided in Appendix~\ref{app:related}.


